%!TEX root = TheoreticalPhysicsC.tex
\section{Lecture: Uncertainty}
We now consider free particles, described in the Hilbert space $\mathcal{H}\equiv L^2(\mathbb{R}^3)$. The space $\mathbb{R}^3$ is interpreted here as position space, that is, in this representation of the Hilbert space the functions of the position observable are represented as multiplication operators. In a pure state $\rho = |\Psi\rangle\langle \Psi|$ the probability density for the outcomes of position measurements is therefore $p(\mathbf{x})\equiv |\Psi(\mathbf{x})|^2$.

A position measurement is a grid of yes/no detectors, essentially a camera; the operator $\widehat{x}_j$ itself is never measured directly -- it stands for this procedure.

The expectation value of the position component $\widehat{x}_j$ is obtained by integration:
\begin{equation}
	\mathbb{E}[x_j] = \iiint x_j |\Psi(\mathbf{x})|^2\, dxdydz = \tr(\rho \widehat{x}_j),
\end{equation}
and correspondingly for the higher moments. In particular, then, the variance of the distribution of $x_j$ is:
\begin{equation}
	\text{var}(x_j) = \tr(\rho \widehat{x}_j^2) - \tr(\rho \widehat{x}_j)^2.
\end{equation}

Analogously, for the momentum operator we obtain
\begin{align}
	\mathbb{E}[p_j] &= \tr(\rho \widehat{p}_j) \\
	\text{var}(p_j) &= \tr(\rho \widehat{p}_j^2) - \tr(\rho \widehat{p}_j)^2
\end{align}

These are two separate experiments: either you measure the position, or you measure the momentum. Up to this point the predictions of classical physics and quantum mechanics do not differ; the break comes now.

It turned out that the two distributions cannot both be sharply concentrated at the same time. That is: the two variances cannot both become arbitrarily small. Let $A$ and $B$ be arbitrary hermitian operators. Then, for every
$\lambda \in \mathbb{R}$, the operator $M = (A+i\lambda B)^\dag (A+i\lambda B)$ is (by definition) positive,
and hence has positive expectation values. That $\lambda$ has to be real is essential: taking the adjoint
conjugates complex numbers, and for complex $\lambda$ the calculation would go wrong. The operator $M$
is a trick; all that matters is that it has the form $C^\dag C$ with $C=A+i\lambda B$, for it is this form that guarantees positivity. Hence, for $\lambda\in \mathbb{R}$ and an arbitrary
density matrix $\rho$:
\begin{equation}
	\tr(\rho M) = \tr(\rho A^2)+\lambda^2\tr(\rho B^2) + \lambda\tr(\rho i[A,B])\ge 0.
\end{equation}
If you insert the value of $\lambda$ that makes the left-hand side smallest (exercise),
you get:
\begin{equation}
	\tr(\rho A^2)\tr(\rho B^2) \ge \frac14 \tr(\rho i[A,B])^2.
\end{equation}
The recipe for this: differentiate with respect to $\lambda$, set the derivative to zero. The inequality holds for arbitrary hermitian operators.
If you now insert $A=\widehat{p}_j - \tr(\rho \widehat{p}_j )\mathbb{I}$ and $B=\widehat{x}_j - \tr(\rho \widehat{x}_j )\mathbb{I}$ and use the canonical commutation relations on the right-hand side, you get
\begin{equation}
	\eqbox{\text{var}(p_j)\text{var}(x_j) \ge \frac{\hbar^2}{4}, \quad\text{or}\quad (\Delta p_j)(\Delta x_j) \ge \frac{\hbar}{2},}
\end{equation}
if you introduce the $(\Delta A)$ as the square root of the variance (the standard deviation). This is the \emph{Heisenberg uncertainty relation}, which in this form is due to Kennard.

The interpretation of the Heisenberg uncertainty relation is this: it is impossible to prepare quantum mechanical particles in such a way that the variances of momentum and position measurements both come out arbitrarily small. The quantities $(\Delta p_j)$ and $(\Delta x_j)$ thus refer to completely different statistical experiments. A standard deviation arises only through many repetitions: prepare, measure, prepare again in exactly the same way, measure again; you determine $(\Delta p_j)$ in a separate run of its own.

% fig_vl09_unschaerfe.tex — Woher die Schranke kommt: die Positivität von
% M = C^dag C ist eine Parabel in lambda, die die Nulllinie nicht
% unterschreiten darf, und genau ihre Scheitelhöhe ist Kennards
% Ungleichung.
%
% Follows the conventions of fig_vl11_kasten.tex.
%
% WARUM NICHT DIE GAUSS-PAKETE.  Die naheliegende Tafel — schmales
% Ortspaket, breites Impulspaket — steht bereits in vl08
% (fig_vl08_fourier): dort sind drei Pakete in beiden Darstellungen
% nebeneinandergestellt und mit sigma_x sigma_p = hbar/2 beschriftet.
% Sie hier zu wiederholen hieße, die vorige Vorlesung noch einmal zu
% zeichnen.  Neu in vl09 ist der *Beweis*: der Kunstgriff mit M, das
% Minimieren über lambda, und was daraus für zulässige Paare
% (Delta x, Delta p) folgt.  Genau das ist hier gezeichnet — Panel (a)
% der Mechanismus, Panel (b) sein Ergebnis.
%
% GEOMETRIE AUS DER FORMEL.  Alles folgt aus der Ungleichung der
% Vorlesung, mit A = p - <p> und B = x - <x>:
%
%       tr(rho M) = tr(rho A^2) + lambda^2 tr(rho B^2)
%                                       + lambda tr(rho i[A,B]) >= 0 .
%
% Hier ist tr(rho A^2) = (Delta p)^2, tr(rho B^2) = (Delta x)^2 und,
% wegen [x,p] = i hbar, tr(rho i[A,B]) = i[p,x] = hbar.  Also
%
%       f(lambda) = (Delta p)^2 + lambda^2 (Delta x)^2 + lambda hbar .
%
% In den dimensionslosen Variablen
%
%       mu    = lambda (Delta x)^2 / hbar ,
%       g(mu) = f (Delta x)^2 / hbar^2 ,
%       kappa = Delta x Delta p / hbar
%
% lautet das
%
%       g(mu) = mu^2 + mu + kappa^2 ,
%
% eine Parabelschar, die NUR vom Produkt kappa abhängt — die einzige
% dimensionslose Größe des Problems.  Scheitel bei mu = -1/2, das heißt
% lambda_* = -hbar / 2 (Delta x)^2, genau der Wert, den die Vorlesung als
% Übung ausrechnen lässt; Scheitelhöhe kappa^2 - 1/4.  Positivität heißt
% also kappa >= 1/2, und das ist Delta x Delta p >= hbar/2.  Die drei
% Kurven sind demnach nicht gewählt, sondern abgelesen: kappa = 1
% (Ungleichung strikt), kappa = 1/2 (Gleichheit, Gauß), kappa = 1/4
% (Scheitel unter null — kann kein Zustand sein).
%
% ENTSCHEIDUNGEN
%   * die y-Achse ist so eng gefasst, dass der Höhenunterschied
%     kappa^2 = 1/4 vs 1/16 (also 0.19) noch gut 5 mm misst: die drei
%     Parabeln sind exakte Vertikalverschiebungen voneinander, weiter
%     lassen sie sich nicht trennen, und die Aussage steckt allein in
%     der Scheitelhöhe;
%   * schraffiert ist in (a) der Halbraum tr(rho M) < 0 und in (b) das
%     Gebiet Delta x Delta p < hbar/2 — zweimal dasselbe Verbot, einmal
%     als Operatoraussage, einmal als Aussage über Messstatistiken;
%   * Farbe kodiert das Produkt kappa und sonst nichts, dieselben drei
%     Farben in beiden Panels: die Parabel in (a) ist als Punkt in (b)
%     wiederzufinden.  Rot und gestrichelt ist der Fall, den es nicht
%     gibt;
%   * `axis on top' in (a): die Schraffur ist unbeschnitten (clip=false,
%     damit die Direktbeschriftung rechts stehen kann) und übermalt sonst
%     die Achslinie, die dann gestrichelt aussieht;
%   * grau ist ausschließlich Referenz: die Nulllinie, die Scheitel-
%     senkrechte lambda_*, und in (b) die Senkrechte Delta x = a, auf der
%     die drei Zustände aus (a) liegen;
%   * die beiden zusätzlichen Punkte auf der Hyperbel in (b) sind
%     dieselbe Gauß-Familie bei anderer Breite: quetscht man den Ort,
%     rutscht der Punkt die Hyperbel hinauf — das ist der Pfeil.  Die
%     Größenordnungsnotiz macht aus der willkürlichen Einheit a eine
%     Aussage über ein Elektron: auf einen Atomdurchmesser eingesperrt,
%     streut seine Geschwindigkeit um einige hundert km/s.
%
% Panel (a) die Parabelschar g(mu) = mu^2 + mu + kappa^2.
% Panel (b) die zulässigen Paare (Delta x, Delta p).

\begin{center}
\begin{tikzpicture}

% ─────────────────────────────────────────────────────────────
%  (a)  Die Positivitätsparabel
% ─────────────────────────────────────────────────────────────
\begin{axis}[
    name=parab,
    tpcplot,
    width=7.5cm, height=6.0cm,
    grid=none,
    xmin=-1.35, xmax=1.15,
    ymin=-0.40, ymax=1.78,
    xtick={-1,-0.5,0,0.5},
    xticklabels={$-1$,$-\frac12$,$0$,$\frac12$},
    ytick={0,0.5,1,1.5},
    xlabel={$\lambda\,(\Delta x)^{2}/\hbar$},
    ylabel={$\tr(\rho M)\,(\Delta x)^{2}/\hbar^{2}$},
    every axis x label/.style={at={(axis description cs:0.5,-0.15)},
                               anchor=north, font=\small},
    every axis y label/.style={at={(axis description cs:-0.14,0.5)},
                               rotate=90, anchor=south, font=\small},
    title={(a) \tpcifenglish{A parabola that may not dip}%
                            {Eine Parabel, die nicht eintauchen darf}},
    samples=181, clip=false, axis on top,
]

  % Verbotener Halbraum: M = C^dag C ist positiv, also tr(rho M) >= 0.
  \path[forbidden] (axis cs:-1.35,-0.40) rectangle (axis cs:1.15,0);
  \draw[levellinegrey] (axis cs:-1.35,0) -- (axis cs:1.15,0);

  % Scheitelsenkrechte mu = -1/2, also lambda_* = -hbar/2(Delta x)^2.
  \draw[classicallim] (axis cs:-0.5,-0.40) -- (axis cs:-0.5,1.10);

  % g(mu) = mu^2 + mu + kappa^2, kappa = Delta x Delta p / hbar.
  \addplot[wavefn, munsell,    domain=-1.35:0.35] {x^2 + x + 1};
  \addplot[wavefn, spacecadet, domain=-1.35:0.35] {x^2 + x + 0.25};
  \addplot[wavefn, birrenred, densely dashed, domain=-1.35:0.35]
    {x^2 + x + 0.0625};

  % Die Scheitel: das ist die ganze Information.
  \addplot[only marks, mark=*, mark size=2.0pt, munsell]
    coordinates {(-0.5,0.75)};
  \addplot[only marks, mark=*, mark size=2.0pt, spacecadet]
    coordinates {(-0.5,0)};
  \addplot[only marks, mark=*, mark size=2.0pt, birrenred]
    coordinates {(-0.5,-0.1875)};

  % Direktbeschriftung am rechten Kurvenende, in der Kurvenfarbe.
  \node[lbl, text=birrengray, font=\scriptsize, anchor=west]
    at (axis cs:0.42,1.66) {$\Delta x\,\Delta p$:};
  \node[lbl, text=munsell,    anchor=west] at (axis cs:0.42,1.4725)
    {$\hbar$};
  \node[lbl, text=spacecadet, anchor=west] at (axis cs:0.42,0.7225)
    {$\hbar/2$};
  \node[lbl, text=birrenred,  anchor=west] at (axis cs:0.42,0.5350)
    {$\hbar/4$};

  % Wo der Scheitel liegt und wie hoch er ist -- ein Block über der
  % Scheitelsenkrechten, im freien Feld über der munsell-Parabel.  Das
  % Minimieren über lambda, das die Vorlesung als Übung stellt, steht
  % damit als Bild da: der Scheitel sinkt mit dem Produkt.
  \node[lbl, text=birrengray, font=\scriptsize, anchor=south, align=center]
    at (axis cs:-0.5,1.12)
    {$\lambda_{*}(\Delta x)^{2}=-\hbar/2$\\[-1pt]
     \tpcifenglish{vertex height}{Scheitelhöhe}\\[-1pt]
     $(\Delta x\,\Delta p/\hbar)^{2}-\tfrac14$};

  % Warum der schraffierte Halbraum leer ist.
  \node[lbl, text=birrenred, font=\scriptsize, anchor=south east,
        fill=white, fill opacity=0.85, text opacity=1, inner sep=2pt]
    at (axis cs:1.13,-0.36)
    {$M=C^{\dag}C\ge 0$: \tpcifenglish{impossible}{unmöglich}};
\end{axis}

% ─────────────────────────────────────────────────────────────
%  (b)  Die zulässigen Paare (Delta x, Delta p)
% ─────────────────────────────────────────────────────────────
\begin{axis}[
    tpcplot,
    at={(parab.south east)}, anchor=south west, xshift=10mm,
    width=6.5cm, height=6.0cm,
    xmin=0, xmax=2.62, ymin=0, ymax=1.85,
    xtick={0,1,2}, ytick={0,0.5,1,1.5},
    grid=major,
    xlabel={$\Delta x/a$},
    ylabel={$\Delta p\,a/\hbar$},
    every axis x label/.style={at={(axis description cs:0.5,-0.15)},
                               anchor=north, font=\small},
    every axis y label/.style={at={(axis description cs:-0.125,0.5)},
                               rotate=90, anchor=south, font=\small},
    title={(b) \tpcifenglish{Which pairs survive}{Welche Paare übrig bleiben}},
    clip=false,
]
  % Verboten: Delta x Delta p < hbar/2.  Links von x = 1/(2 ymax) hat die
  % Hyperbel den Rahmen schon verlassen, dort wird der Streifen ganz
  % schraffiert.
  \path[forbidden] (axis cs:0,0) rectangle (axis cs:0.27027,1.85);
  \addplot[forbidden, draw=none, domain=0.27027:2.62, samples=200]
    {0.5/x} \closedcycle;

  % Die Schranke selbst.
  \addplot[spacecadet, line width=1.1pt, domain=0.27027:2.62, samples=200]
    {0.5/x};

  % Die Senkrechte Delta x = a trägt die drei Zustände aus (a).
  \draw[classicallim] (axis cs:1,0) -- (axis cs:1,1.02);

  \addplot[only marks, mark=*, mark size=2.2pt, munsell]
    coordinates {(1,1)};
  \addplot[only marks, mark=*, mark size=2.2pt, spacecadet]
    coordinates {(1,0.5)};
  % Ring thickness goes on the plot itself: nested inside mark options it
  % is typeset rather than applied, which shows up as nullfont warnings.
  \addplot[only marks, mark=o, mark size=2.6pt, birrenred, line width=1.0pt]
    coordinates {(1,0.25)};

  % Dieselbe Gauß-Familie bei anderer Breite: quetschen heißt die
  % Hyperbel entlangrutschen.
  \addplot[only marks, mark=*, mark size=2.2pt, spacecadet]
    coordinates {(0.5,1) (2,0.25)};
  \draw[-{Stealth[length=4pt]}, spacecadet, line width=0.7pt]
    (axis cs:1.05,0.60) .. controls (axis cs:0.76,0.74) ..
    (axis cs:0.60,0.96);

  \node[lbl, text=spacecadet, font=\scriptsize, anchor=south west]
    at (axis cs:0.58,1.08) {\tpcifenglish{squeeze $x$}{$x$ quetschen}};
  \node[lbl, anchor=south east, align=right] at (axis cs:2.58,0.40)
    {\tpcifenglish{Gaussian: equality}{Gauß: Gleichheit}\\[-1pt]
     $\Delta x\,\Delta p=\hbar/2$};

  % Das Verbot beschriftet den Streifen, in dem es gilt, statt die Mitte
  % der Tafel zu belegen.
  \node[lbl, text=birrengray, font=\scriptsize, rotate=90, anchor=center]
    at (axis cs:0.135,1.00) {\tpcifenglish{no state}{kein Zustand}};

  % Die Brücke zu Panel (a): gleiche Farben, gleiche drei Zustände.
  \node[lbl, text=birrengray, font=\scriptsize, anchor=west]
    at (axis cs:1.07,0.87) {\tpcifenglish{the three of (a)}{die drei aus (a)}};

  % Größenordnung, damit die Einheit a etwas bedeutet.
  \node[lbl, text=birrengray, font=\scriptsize, anchor=north east,
        align=right] at (axis cs:2.60,1.82)
    {$e^{-}$, $a=10^{-10}\,$m:\\[-1pt]
     $\Delta v\gtrsim 6\cdot10^{5}\,$m/s};
\end{axis}

\end{tikzpicture}
\end{center}


Panel~(a) shows what the minimisation over $\lambda$ does: if the product
$\Delta x\,\Delta p$ falls, the vertex falls with it, and at $\hbar/2$ it just
touches zero. Anything below that would have a negative expectation value.
Panel~(b) translates the same prohibition into measured quantities: only what lies on
or above the hyperbola $\Delta x\,\Delta p=\hbar/2$ is allowed, and whoever squeezes the
position slides up along it.

\begin{merksatz}[What the uncertainty relation does not say]
It says nothing about successive measurements. The widespread account, in which you first measure the position and then the momentum on the same particle, is wrong, even though nearly all authors present it that way. For sequences of measurements there is an inequality as well, but a different and more complicated one. This distinction matters for the examination.
\end{merksatz}